\section{追う量を変える}\label{節：追う量を変える}
四つの基本型を目標に,与えられた数列から追う量を作ろう.

  \subsection{補助数列を作る}\label{節：補助数列の設計}
    元の\bookterm{sequence}{数列}から作った,更新を調べるための別の\bookterm{sequence}{数列}を\textbf{\bookterm{auxiliary}{補助数列}}という.
    先に$b_{n+1}=pb_n$や$b_{n+1}-b_n=f(n)$という目標を決め,必要な条件を元の式から求める.
    大切なのは,\textbf{何を消し,どの更新を残すか}である.

    \bookstep{引いて付加項を消す}
    \sectionref*{節：漸化式を解く前に}の例$a_{n+1}=2a_n+3$で$b_n=a_n+c$とおくと,
    \[
      b_{n+1}=2b_n+3-c
    \]
    である.等比型にしたいので$3-c=0$を要求すると,$c=3$が決まる.
    $a_n+3=a_n-(-3)$と見れば,定数解$-3$を引いたことになる.初項によらず,この量は$2$倍される.

    一般に$a_{n+1}=pa_n+f(n)$なら,同じ更新をする既知の量$g(n)$,
    すなわち
    \[
      g(n+1)=pg(n)+f(n)
    \]
    を探す.必要なのは\textbf{一つの解}であり,元の初項まで合わせる必要はない.
    $b_n=a_n-g(n)$とすれば,二つの式を引いて
    \[
      b_{n+1}=pb_n
    \]
    となる.定数項には定数,一次式には一次式を候補にして係数を比較し,条件を満たせなければ選び直す.\footnote{候補の次数や係数の決め方は\sectionref*{節：階差等比複合型}で扱う.}
    ただし,その量を簡単に求められるとは限らない.

    \bookstep{割って倍率をそろえる}
    例えば$a_{n+1}=\frac{n+1}{n}a_n+(n+1)$では,両辺を$n+1$で割ると,
    \[
      \frac{a_{n+1}}{n+1}=\frac{a_n}{n}+1
    \]
    となる.$b_n=\frac{a_n}{n}$を追えば,更新は$b_{n+1}=b_n+1$である.
    $b_1=a_1$なので,$b_n=a_1+n-1$から$a_n=n(a_1+n-1)$へ戻せる.
    割る量の倍率で,元の倍率を打ち消したのである.

    一般に$a_{n+1}=p_na_n+q_n$で倍率$p_n$が添字ごとに変わるなら,
    $p_n\ne0$の範囲で,$h_{n+1}=p_nh_n$を満たす非零の\bookterm{sequence}{数列}$h_n$を選ぶ.
    $h_1=1$から倍率を掛け合わせれば作れる.先ほどは$h_n=n$であった.
    $b_n=\frac{a_n}{h_n}$と置くと,
    \[
      b_{n+1}-b_n=\frac{q_n}{h_{n+1}}
    \]
    となり,残った変化を足せばよい.\footnote{具体的には$h_n=\prod_{k=1}^{n-1}p_k$と選べる.ただし,その積を簡単に計算できるかは別に調べる.
    このように基準をそろえる操作を\textbf{\bookterm{normalize}{正規化}}という.
    具体例と既知解を引く別解は\sectionref*{節：変数係数の一次漸化式}で扱う.
    \sectionref*{節：階比型}で因子を掛けた方法とは,$g(n)=1/h_n$と対応する.割る量と掛ける量は互いに逆数である.}
    同じ式にも引く方法と割る方法を使える.条件と残る計算を比べて選ぼう.

    \bookstep{項をまとめる}
    直前の二項で次が決まるなら,二項を合わせた量も候補になる.
    例えば$a_{n+2}=5a_{n+1}-6a_n$では,
    \[
      a_{n+2}-2a_{n+1}=3(a_{n+1}-2a_n)
    \]
    とまとめると,$b_n=a_{n+1}-2a_n$が$b_{n+1}=3b_n$を満たす.
    従って,$b_n=(a_2-2a_1)3^{n-1}$が分かる.
    元へ戻すには,
    \[
      a_{n+1}=2a_n+b_n
    \]
    を初項$a_1$とともに解く必要がある.
    この復元は,\bookproblemref{example-three-term-1}{1}{隣接3項型の例題(1)}で最後まで計算する.

    引く係数$2$や残る倍率$3$は,$b_n=a_{n+1}-\alpha a_n$と置き,
    $b_{n+1}=\beta b_n$を目指して係数を合わせれば見つかる.
    左辺に元の式を入れると$(5-\alpha)a_{n+1}-6a_n$,右辺は$\beta a_{n+1}-\alpha\beta a_n$になる.
    この例なら$\alpha+\beta=5,\ \alpha\beta=6$が条件になる.
    目標とする更新から,係数が決まった.\footnote{この考え方は\sectionref*{節：隣接3項間漸化式型}や\sectionref*{節：連立型}へつながる.}

    どの置換でも,\bookterm{initial}{初期条件}と復元を確かめる.
    $a_n=b_n+g(n)$や$a_n=h_nb_n$のように戻し,複数の項をまとめた場合は,その関係から元の\bookterm{sequence}{数列}を解く.
    逆数・\bookterm{logarithm}{対数}などを定義できる条件も必要である.

  \subsection{発見と証明}

  消す項やそろえる倍率を選ぶのが\textbf{発想},代入して\bookterm{general}{一般項}を求め,元へ戻すのが\textbf{計算}である.
  逆数・\bookterm{logarithm}{対数}・割り算などの条件は,必要になった時点で確認する.

  数項を計算する実験は候補を探す助けになるが,数項の一致は証明ではない.変形や実験から得た式は,帰納法でも確かめられる.\footnote{予想した式が\bookterm{initial}{初期条件}と元の更新式を満たせば,順に次の項が一意に決まる\bookterm{recurrence}{漸化式}では,\bookterm{induction}{数学的帰納法}で一致を示せる.
  \sectionref*{節：数学的帰納法}を参照してほしい.}

  例えば,\sectionref*{節：数列と漸化式}の$a_1=1,\ a_{n+1}=a_n+2$で$a_n=2n-1$を確かめよう.
  $n=1$では$2\cdot1-1=a_1$である.$a_k=2k-1$を仮定すると,
  \[
    a_{k+1}=a_k+2=(2k-1)+2=2(k+1)-1
  \]
  となる.出発点と一つ先への移行から,すべての$n\ge1$で正しいと示す方法が\mbox{\bookterm{induction}{数学的帰納法}}である.

  \bookterm{general}{一般項}を簡単な式で表すことに加えて,正値・\bookterm{limit}{極限}・整数性などを調べる目標もある.
  \chapterref{章：実際の解き方}では方法の選択を練習し,発展演習と付録では目標の広がりを見る.
